\section{Introduction}

Large language models (LLMs) are not only language technologies; they are also increasingly being studied as objects of behavioural experimentation. Researchers have used persona-conditioned LLM responses to simulate respondents with different characteristics \cite{Argyle2023OutOfOneMany}, attempted to reproduce experimental effects previously observed in human participants \cite{Aher2023SimulateHumans}, and treated models as simulated economic agents whose decisions can be examined experimentally \cite{Horton2023HomoSilicus}. Cognitive tasks have likewise been used to characterise behavioural regularities and reasoning biases across GPT models \cite{Binz2023CognitivePsychologyGPT3,Hagendorff2023HumanLikeBiases}. LLMs therefore offer a potentially useful methodological resource: their behaviour can be tested repeatedly under controlled task manipulations.

Nevertheless, a single LLM run is not equivalent to a human participant, and behavioural similarity does not establish a shared cognitive mechanism \cite{Shanahan2023RolePlay}. An LLM output is a conditional sample generated in response to a particular input under specified generation settings. The scientific value of studying an LLM as a cognitive model therefore depends not only on the task-specific behaviour observed, or on its proximity to a human reference sample, but also on whether the resulting conclusions are robust to predefined variations in prompt formulation \cite{Loya2023PromptSensitivityDecision}. In this study, the term \emph{cognitive model} denotes a computational model capable of producing comparable behaviour in cognitive tasks; it does not imply that an LLM possesses the same underlying cognitive mechanisms as humans.

Prompt formulation is central to this question of reliability. A prompt does more than communicate task rules: it can assign the model a role, make particular information salient, and constrain the form of its response. Variations in prompt formulation and generation temperature can alter model performance or choice behaviour \cite{Loya2023PromptSensitivityDecision}. Moreover, framing a model as an experimental participant may elicit a particular simulated identity rather than reveal a stable, context-independent disposition \cite{Shanahan2023RolePlay}. Prompts should therefore be treated as part of the measurement procedure in behavioural research. If prompts that preserve the same task information and rules produce different patterns of exploration, learning, or risk-taking, the resulting conclusions may partly reflect the measurement procedure rather than stable features of the model's behaviour in the target task.

Although previous studies have established that prompting can affect model outputs, the reliability of LLM behaviour has not yet been fully characterised. Many evaluations focus on accuracy, aggregate human similarity, or individual empirical effects, while studies involving multiple prompt templates have primarily examined formatting variations in static benchmarks \cite{Sclar2024PromptFormatting}. Loya et al.~\cite{Loya2023PromptSensitivityDecision} extended this line of inquiry to the sequential Horizon Task, but considered only one decision paradigm. Relatively few studies combine full task runs as the unit of replication; predefined prompt variants that preserve task rules and available information; task- and metric-specific comparisons of model-by-prompt variation under a common implementation; and separate assessments of prompt stability, task-specific behavioural outcomes, and human-reference proximity. Without this combination, it is difficult to determine whether an observed behavioural conclusion reflects a stable feature of a model's task behaviour or a consequence of a particular measurement formulation. These dimensions are not interchangeable: a model may be stable but dissimilar to humans, resemble humans only under a particular prompt, or vary across prompts while maintaining similar average rewards.

Language constitutes a further source of contextual variation. Multilingual evaluations indicate that the performance of a single model may vary substantially across languages and tasks \cite{Ahuja2023MEGA}; evidence obtained in English therefore cannot automatically be generalised to other languages. Interpretable cross-language behavioural comparisons require the model and task environment to be held constant, while task-relevant information and the intended prompt manipulation are aligned across languages. Such alignment does not, however, imply complete pragmatic equivalence among expressions in different natural languages.

This study uses three sequential decision-making tasks. The Horizon Task examines information-seeking choices and changes in exploration across decision horizons \cite{Wilson2014}. The Iowa Gambling Task (IGT) examines longer-term adaptation to reward--loss feedback \cite{Bechara1994}. The Balloon Analogue Risk Task (BART) examines sequential risk-taking in which potential rewards accumulate but the temporary earnings associated with a balloon may be lost if it bursts \cite{Lejuez2002}. Together, these tasks cover exploration, feedback-based learning, and risk management. Participant-level human data are available for each task, permitting descriptive comparisons between human participants and LLM run-level summaries on aligned behavioural metrics \cite{Feng2021ExploreExploit,Steingroever2015IGTData,Sebri2023BART}. Human-reference proximity is used solely as a descriptive behavioural benchmark and not as evidence of shared cognitive mechanisms.

Across the Horizon Task, IGT, and BART, the study compares a Neutral baseline with three manipulated prompt conditions matched on task rules. The conditions preserve the available actions, feedback structure, and boundaries of hidden information, but are not assumed to be pragmatically identical. In particular, the Task-specific construct emphasis condition deliberately changes the salience of concepts relevant to the task. The English-language comparison evaluates prompt sensitivity in GPT-4.1, GPT-5.4, and GPT-5.4 Mini. The cross-language comparison holds the model fixed at GPT-4.1 and compares English, Simplified Chinese, and Spanish. The replication structure and controls applied to the task environments are described in the Method section.

This study makes four principal contributions. First, it defines behavioural robustness as the stability of behaviour across predefined, task-rule-matched prompt conditions, and formulation sensitivity as the magnitude of change relative to the Neutral baseline. Second, it tests, separately for each task and metric, whether prompt effects vary across models and then describes any recurring patterns across tasks. Third, it reports prompt sensitivity, task-specific behavioural outcomes, and human-reference proximity as distinct evaluation dimensions. Fourth, in a multilingual design that holds the model fixed, it locates the baseline behaviour and within-language prompt effects of English, Simplified Chinese, and Spanish jointly within a common centred coordinate system. Because construct emphasis constitutes a substantive manipulation of attentional focus, an effect in that condition is not automatically interpreted as unreliability arising from incidental wording.

The study addresses the following research questions:

\noindent\textbf{RQ1 --- Prompt sensitivity.}
To what extent do Instruction specificity, Role framing, and Task-specific construct emphasis alter a model's task-specific behaviour relative to the Neutral baseline?

\noindent\textbf{RQ2 --- Cross-model variation.}
Within each task and behavioural metric, do the prompt effects of GPT-4.1 or GPT-5.4 Mini differ from those of GPT-5.4? Do the directions and magnitudes of these differences exhibit recurring descriptive patterns across tasks and metrics?

\noindent\textbf{RQ3 --- Cross-language variation.}
When the model, task environment, and generation parameters are held constant, and the intended prompt manipulations and task-relevant information are semantically aligned across languages, do English, Simplified Chinese, and Spanish produce different baseline behaviours or within-language prompt effects?

\noindent\textbf{RQ4 --- Prompt-associated changes in human-reference proximity.}
Relative to the Neutral baseline, do alternative prompt formulations change the absolute human-SD-scaled mean deviation from the human reference mean or the proportion of runs falling within the empirical human reference interval?

RQ4 is explicitly descriptive. Here, ``relative to the Neutral baseline'' refers to comparing the position of each prompt condition with that of the Neutral condition against the same human reference distribution; it does not involve combining Neutral runs with human participants. The corresponding analyses report the signed and absolute human-SD-scaled mean deviations and empirical reference-interval coverage. These quantities are not treated as Cohen's \(d\), are not subjected to binary significance decisions, and are not interpreted as evidence of shared cognitive mechanisms.

Across the model snapshots, tasks, and language implementations examined here, prompt formulations were associated with changes in several task-specific behaviours, but the direction and magnitude of these changes varied by model, task, metric, and language condition. No formulation produced an effect of consistent direction and magnitude across all settings. Direct cross-model interaction contrasts and symmetric, jointly centred three-language comparisons further showed that behavioural conclusions obtained from a single model or a prompt in a single language cannot be assumed to generalise to other models or language conditions. The descriptive human-reference analyses also indicated that prompt formulation could change both the position of the model mean relative to the human mean and the coverage of the empirical human interval by run-level values; these two quantities did not invariably move together. Overall, the robustness of LLM behavioural evidence is conditional on the complete measurement specification, rather than an intrinsic model property separable from the prompt, task, language, and outcome metric.

The remainder of this dissertation is organised as follows. Background and Related Work synthesises research on LLM behaviour, prompt sensitivity, cross-model and cross-language variation, sequential decision tasks, and human-reference comparisons, and uses this literature to define the research gap. Method describes the study design, tasks and prompt materials, experimental procedure, data-quality controls, human reference data, behavioural metrics, statistical analyses, and reproducibility boundaries. Results first reports the analysis sample and data quality, then addresses the four research questions through within-model prompt effects, cross-model contrasts, joint three-language centred comparisons, and prompt-associated changes in human-reference proximity. Discussion integrates these findings by considering prompt formulation as part of the behavioural measurement procedure, the methodological contribution of the study, the interpretive limits of human-reference comparisons, and the implications of the study's limitations for replication. Finally, Conclusion and Future Work answers the central research objective, summarises the principal findings and contributions, and identifies priorities for subsequent work.
